%! Inverse Trigonometric Functions — Unit 3, Lesson 3.9 — Homework (Answer Key)
\documentclass[10pt]{article}
\usepackage{precalculus-article}
\usepackage{precalculus-key}   % pulls in -boxes; do NOT also load -boxes
\newcommand{\TallMath}[1]{$\displaystyle #1\rule[-1.4em]{0pt}{3.2em}$}

\begin{document}

\pageheader{Unit 3, Lesson 3.9}{Homework --- Answer Key}

\vspace{0.10in}

\begin{notesbox}{Part 1 --- Evaluate --- Key}

\renewcommand{\arraystretch}{2.0}
\begin{tabularx}{\linewidth}{@{} X X X @{}}
  \textbf{1.} $\arcsin\!\left(\dfrac{\sqrt{2}}{2}\right) = $ \ans{$\dfrac{\pi}{4}$}
  &
  \textbf{2.} $\arccos\!\left(-\dfrac{1}{2}\right) = $ \ans{$\dfrac{2\pi}{3}$}
  &
  \textbf{3.} $\arctan\!\left(\dfrac{\sqrt{3}}{3}\right) = $ \ans{$\dfrac{\pi}{6}$} \\[4pt]
  \textbf{4.} $\sin^{-1}(0) = $ \ans{$0$}
  &
  \textbf{5.} $\cos^{-1}(1) = $ \ans{$0$}
  &
  \textbf{6.} $\tan^{-1}\!\left(-\dfrac{\sqrt{3}}{3}\right) = $ \ans{$-\dfrac{\pi}{6}$}
\end{tabularx}

\end{notesbox}

\vspace{0.08in}

\begin{notesbox}{Part 2 --- Domain and Range --- Key}

\textbf{7.}

\renewcommand{\arraystretch}{1.9}
\begin{tabular}{|l|c|c|}
\hline
\textbf{Function} & \textbf{Domain} & \textbf{Range} \\
\hline
$y = \arcsin x$    &
  \ans{$[-1,\,1]$} &
  \ans{$\left[-\dfrac{\pi}{2},\,\dfrac{\pi}{2}\right]$} \\
\hline
$y = \arccos x$    &
  \ans{$[-1,\,1]$} &
  \ans{$[0,\,\pi]$} \\
\hline
$y = \arctan x$    &
  \ans{$(-\infty,\infty)$} &
  \ans{$\left(-\dfrac{\pi}{2},\,\dfrac{\pi}{2}\right)$} \\
\hline
\end{tabular}

\end{notesbox}

\vspace{0.08in}

\begin{notesbox}{Part 3 --- AP-Style Multiple Choice --- Key}

\textbf{8.} \ans{(D) $\dfrac{3\pi}{4}$}

$\cos\!\bigl(\tfrac{3\pi}{4}\bigr) = -\dfrac{\sqrt{2}}{2}$ and
$\tfrac{3\pi}{4} \in [0,\pi]$ (restricted domain of cosine).

\vspace{0.08in}

\textbf{9.} \ans{(C) The range of $f$ is $\left(-\dfrac{\pi}{2},\,\dfrac{\pi}{2}\right)$.}

The range of $\arctan$ is the open interval $(-\tfrac{\pi}{2},\tfrac{\pi}{2})$.
(A) is false --- range is not all reals.
(B) is false --- domain of $\arctan$ is all reals.
(D) is false --- $\arctan$ is an odd function, so $f(-x) = -f(x)$ is true.

\end{notesbox}

\vspace{0.08in}

\begin{extensionbox}{Extension --- Composition of Trig and Inverse Trig --- Key}

\textbf{10.} Let $\theta = \arctan x$ with $x > 0$, so $\tan\theta = x$.
Right triangle: opposite $= x$, adjacent $= 1$, hypotenuse $= \sqrt{x^2+1}$.

\ans{$\sin(\arctan x) = \dfrac{x}{\sqrt{x^2+1}}$}

\vspace{0.06in}

\textbf{Preview Lesson 3.10:} $\cos\theta = \tfrac{1}{2}$ has
\ans{two} solutions on $[0,2\pi)$: \ans{$\theta = \dfrac{\pi}{3}$} and
\ans{$\theta = \dfrac{5\pi}{3}$}.
Using $\arccos$: $\arccos\!\bigl(\tfrac{1}{2}\bigr) = \dfrac{\pi}{3}$ gives one solution;
the other is found by symmetry.

\end{extensionbox}

\end{document}
